\naslov{Sturm-Liouvilleova teorija}

Od tu naprej SL teorija.
Naj bo $I = [a,b]$ in $p,q,r: I \to \R$ funkcije, ki zadoščajo
\begin{itemize}
\item $p \in \zvodvi{1}$
\item $q,r \in \zveznei{}$
\item $p,r$ nenegativni
\end{itemize}
Na $\zvodv{2}{I}$ definiramo pridruženi \pojem{Sturm-Liouvilleov} operator
\[
  Lu = (p u')' + qu.
\]
Naj bodo $\alpha_0, \alpha_1, \beta_0, \beta_1 \in \R$ take, da je
$\abs{\alpha_0} + \abs{\alpha_1} > 0$ in $\abs{\beta_0} + \abs{\beta_1} > 0$
(torej paroma ne obe $0$).
Označimo $B_a(u) = \alpha_0 u(a) + \alpha_1 u'(a)$ in $B_b(u) = \beta_0 u(b) +
\beta_1 u'(b)$.
SL problem sprašuje po netrivialnih rešitvah enačbe
\begin{equation}
  \label{eq:sl}
  Lu = - \lambda r u
\end{equation}
ob pogojih $B_a(u) = 0$ in $B_b(u) = 0$.
Funkcijo $r$ imenujemo \pojem{utež} ali \pojem{gostota}.
Problem je \pojem{regularen}, če sta $p, r > 0$, in \pojem{singularen}, če sta
$p$ ali $q$ enaka $0$ v točki $a$ ali $b$, ali če imata v kakšni od teh točk
skok, ali pa če je $I$ neomejen.
Če par $(\lambda, u)$ reši SL problem, je $\lambda$ \pojem{lastna vrednost}, $u$
pa \pojem{lastna funkcija}.

\vprasanje{Definiraj Sturm-Liouvilleov problem. Kdaj je regularen in kdaj
  singularen?}

\begin{trditev}
  Vsako homogeno regularno linearno NDE drugega reda lahko zapišemo v obliki SL
  problema.
\end{trditev}

\begin{proof}
  Imejmo enačbo $A v'' + B v' + C v = 0$.
  Velja $Lv = p v'' + p' v' + (\ldots)$, kjer smo člene ničtega reda skrili.
  Veljati bi moralo $B = A'$.
  Če to ni res, lahko zaradi homogenosti enačbo množimo z neničelno funkcijo $m$
  in dobimo $m(A v'' + B v') + m C v = 0$, kjer $m$ določimo tako, da bo veljalo
  $(mA)' = mB$; izračun pokaže
  \[
	m = \inv{A} \exp\!\left( \int \frac{B}{A} \right).
  \]
  Na koncu dobimo
  \[
	m A v'' + m B v' + m C v = (p v')' + (m C v) = 0
  \]
  za $p = \exp(B/A)$ in $q + \lambda r = \frac{C}{A} p$.
\end{proof}

\vprasanje{Pokaži, da se vsako enačbo drugega reda da zapisati kot SL problem.}

Linearen operator na Banachovem prostoru je zvezen natanko tedaj, ko je omejen.
Ker diferencialni operatorji običajno niso omejeni, je to sila neprijetno, zato
se omejimo na podprostore.
SL problem gledamo na definicijskem območju $L$, ki bo v prihodnje bodisi
\[
  \mathcal{U} = \{u \in \zvodv{2}{I} \such B_a(u) = B_b(u) = 0\}
\]
bodisi
\[
  \mathcal{V} = \{u \in \zvodv{2}{I} \such u(a) = u(b), u'(a) = u'(b)\},
\]
odvisno če gledamo običajen ali periodičen problem (druga množica so ravno
pogoji periodičnega problema).
Operatorja $L_{\mathcal{U}}$ in $L_{\mathcal{V}}$ (torej zožitvi $L$ na ta
prostora) sta definirana na gosti podmnožici $l^2(I)$.
Za zvezno funkcijo $r: I \to \R$ označimo
\[
  l^2(r) = \{f: I \to \C \such f\;\text{merljiva},
  \norm{f}_{l^2(r)}^2 = \int_I \abs{f}^2 r < \infty \}.
\]
Ker je $r$ zvezna na kompaktu, obstajata $c_1$ in $c_2$, da velja $c_1 \le r(x)
\le c_2$, kar označimo z $r \sim 1$.

\begin{opomba}
  Oznako lahko razširimo: $r \sim g$ natanko tedaj, ko obstajata konstanti $c_1$
  in $c_2$, da velja $c_1 g(x) \le r(x) \le c_2 r(x)$.
\end{opomba}

Vidimo, da je $\norm{f}_{l^2(r)} \sim \norm{f}_{l^2}$.

\begin{trditev}
  Naj bo $Lu = (pu')' + qu$ standardni SL operator.
  Če je $v \in \zvodv{2}{I}$, velja $u L v - v L u = [p \cdot (u v' - u' v)]'$
  in
  \[
	\int_I u L v - v L u = \left. p \cdot (u v' - u' v) \right|_{\partial I}.
  \]
\end{trditev}

Dokaz je račun.

\begin{posledica}
  Operator $L_{\mathcal{U}}$ je simetričen.
  Enako velja za $\mathcal{V}$, če je $p(a) = p(b)$.
\end{posledica}

\begin{proof}
  Naj bosta $u$ in $v$ realni funkciji.
  Velja
  \[
	\sk{Lu, v} - \sk{u, Lv}
	= \int_I (vLu - uLv) = \left. p(u' v - u v') \right|_{\partial I}.
  \]
  Če je $p(a) = p(b)$ in $L = L_{\mathcal{V}}$, je to enako $0$.
  V drugem primeru prepišemo robni pogoj v
  \[
	\begin{bmatrix}
	  u(a) & u'(a) \\
	  v(a) & v'(a)
	\end{bmatrix}
	\cdot
	\begin{bmatrix}
	  \alpha_0 \\ \alpha_1
	\end{bmatrix}
	= 0.
  \]
  Privzeli smo, da nista oba $\alpha_0, \alpha_1$ enaka $0$, torej je matrika
  singularna.
  Enako velja za $b$, iz česar sledi $\left. u'v - uv' \right|_a = 0$ in $\left.
	u'v - u v' \right|_b = 0$.

  Če sta $u$ in $v$ kompleksni funkciji, torej $u = u_1 + i u_2$ in $v = v_1 + i
  v_2$, pa velja
  \[
	\sk{Lu, v} = \sk{Lu_1, v_1} + \sk{L u_2, v_2} -i (\sk{Lu_1, v_2} + \sk{L
	  u_2, v_1}).
  \]
  Sedaj le upoštevamo realni primer, dobimo $\sk{Lu, v} = \sk{u, Lv}$.
\end{proof}

\vprasanje{Dokaži: če je $\mathcal{U}$ množica $\zvodv{2}{I}$ funkcij, ki
  ustrezajo robnemu pogoju SL problema, je zožitev operatorja $L_{\mathcal{U}}$
  simetrična.}

\begin{trditev}
  Vse lastne vrednosti operatorjev $L_\cdot$ glede na utež $r$ so realne.
\end{trditev}

\begin{proof}
  Če je $u$ lastna funkcija, velja
  \[
	0 = \sk{Lu, u} - \sk{u, Lu}
	= \sk{- \lambda r u, u} - \sk{u, - \lambda r u}
	= (\konj{\lambda} - \lambda) \sk{u, ru}.
  \]
\end{proof}

\vprasanje{Dokaži: vse lastne vrednosti SL problema so realne.}

\begin{lema}
  Če so funkcije $f_1, \ldots, f_n$ linearno odvisne, je determinanta Wronskega
  identično enaka $0$.
  Obratno, če je $\Lambda$ diferencialen operator z realnimi koeficienti in
  $y_1, \ldots, y_n$ rešijo enačbo $\Lambda y = 0$, ter je $W(y_1, \ldots, y_n)
  = 0$, tedaj so linearno odvisne.
\end{lema}

\begin{trditev}
  Naj bo $\lambda$ lastna vrednost za SL problem $L_{\mathcal{U}} = - \lambda r
  u$.
  Tedaj je pripadajoči lastni podprostor enodimenzionalen.
\end{trditev}

\begin{proof}
  Naj za $u, v \in \mathcal{U}$ velja $Lu = - \lambda r u$ in $Lv = - \lambda r
  v$.
  Sledi $u L v - v L u = 0$, a hkrati $u L v - v L u = [p(u' v - u v')]'$, torej
  je $f = p\,(u'v - uv')$ konstantna.
  Ker sta $u, v \in \mathcal{U}$, je $f(a) = f(b) = 0$ in zato $f = 0$.
  Ker pa je $p \ne 0$, sledi še $u'v - vu' = 0$, kar je ravno determinanta
  Wronskega.
  Iz leme sledi, da sta $u$ in $v$ linearno odvisni, saj obe ležita v jedru
  diferencialnega operatorja $\Lambda y = Ly + \lambda ry$.
\end{proof}

\vprasanje{Dokaži: lastni podprostor za SL problem je enodimenzionalen.}

\begin{izrek}[spektralni izrek za SL problem]
  Za vsak regularen SL problem obstaja zaporedje lastnih vrednosti $\lambda_1 <
  \lambda_2 < \ldots < \infty$, da velja
  \begin{itemize}
  \item $\lim \lambda_n = \infty$
  \item Pripadajoč normaliziran sistem lastnih funkcij $\{v_n\}_n$ sestavlja
	kompleten ortonormiran sistem v $l^2(r)$.
  \end{itemize}
  Enako velja za periodičen SL problem, pri čemer imamo lahko dvodimenzionalne
  lastne podprostore.
\end{izrek}

\vprasanje{Povej spektralni izrek za SL problem.}

% LocalWords:  Sturm-Liouvilleova SL Sturm-Liouvilleov ničtega Banachovem
% LocalWords:  Wronskega
